\section{角} 
\subsection{在 \tkzname{tkz-euclide} 中的定义与用法}
在欧几里得几何中，角是由两条射线（称为角的边）共享一个公共端点（称为角的顶点）所形成的图形 [Wikipedia]。在 \tkzname{tkz-euclide} 中，一条射线由两个点确定，因此角由三个点确定，如 $\widehat{AOB}$；顶点 $O$ 是第二个点。它们的顺序很重要，因为角默认按逆时针方向给出。
在三角学以及一般数学中，平面角约定以逆时针方向度量，从指向正右方（或东方）的 $0^\circ$ 开始，到指向正上方（或北方）的 $90^\circ$ [Wikipedia]。
我们约定：逆时针度量的角为正。

  \begin{center}
    \begin{tikzpicture}[scale=.75]
      \node{顺时针};
      \tkzDefPoint(0,0){O} \tkzDefPoint(90:2){A}\tkzDefPoint(180:2){B}   
      \tkzDrawArc[black,line width=2pt,arrows = {Stealth-}](O,B)(A)
    \end{tikzpicture}
    \begin{tikzpicture}[scale=.75]
          \node{逆时针};
      \tkzDefPoint(0,0){O} \tkzDefPoint(90:2){A}\tkzDefPoint(0:2){B}   
      \tkzDrawArc[black,line width=2pt,arrows = {-Stealth}](O,A)(B)
    \end{tikzpicture}  
  \end{center}
  
 角涉及多个宏命令，例如 \tkzcname{tkzDefPoint}、\tkzcname{tkzDefPointBy[rotation = \dots]}、\tkzcname{tkzDrawArc}
 以及下文中介绍的 \tkzcname{tkzGetAngle}。除最后一个外，所有这些宏命令都接受负角。
 
 \begin{figure}[!ht]
 \centering
 \begin{tabular}{|c|c|}
 \hline
 \tkzsubf{\begin{tikzpicture}
 \tkzDefPoint(0,0){O}    \tkzDefPoint(0:2){A}
 \tkzDefPointBy[rotation=center O angle 80](A)  \tkzGetPoint{B}
 \tkzDrawSegments[-{Stealth}](O,A O,B)
 \tkzMarkAngles[size=2,-Stealth,teal](A,O,B)
 \tkzFindAngle(A,O,B)   \tkzGetAngle{an}
 \tkzLabelAngle[pos=1,teal](A,O,B){$ \pgfmathprintnumber{\an}^\circ$}
 \tkzLabelPoints(A)  \tkzLabelPoints[above](B)
 \end{tikzpicture}}
      {从 $(O,A)$ 旋转 $80^\circ$ 到 $(O,B)$\\ 
    {\textbackslash}tkzDefPointBy[rotation=center O angle 80]}
 &
 \tkzsubf{\begin{tikzpicture}
 \tkzDefPoint(0,0){O}    \tkzDefPoint(0:2){A}
 \tkzDefPointBy[rotation=center O angle -80](A)  \tkzGetPoint{B}
 \tkzDrawSegments[-{Stealth}](O,A O,B)
 \tkzMarkAngles[size=2,Stealth-,red](B,O,A)
 \tkzFindAngle(B,O,A)   \tkzGetAngle{an}
 \tkzLabelAngle[pos=1,red](B,O,A){$-\pgfmathprintnumber{\an}^\circ$}
\tkzLabelPoints[right](A)  \tkzLabelPoints[below](B)
 \end{tikzpicture}}
  {从 $(O,A)$ 旋转 $-80^\circ$ 到 $(O,B)$\\ 
     {\textbackslash}tkzDefPointBy[rotation=center O angle -80]}
 \\ \hline
 \tkzsubf{\begin{tikzpicture}
 \tkzDefPoint(0,0){O}    \tkzDefPoint(0:2){A}
 \tkzDefPointBy[rotation=center O angle 80](A)  \tkzGetPoint{B}
 \tkzDrawSegments[-{Stealth}](O,A O,B)
 \tkzMarkAngles[size=1.5,-Stealth,teal](A,O,B)
 \tkzFindAngle(A,O,B)   \tkzGetAngle{an}
 \tkzLabelAngle[pos=1,teal](A,O,B){$ \pgfmathprintnumber{\an}^\circ$}
\tkzLabelPoints(A)  \tkzLabelPoints[above](B)
 \end{tikzpicture}}
      { {\textbackslash}tkzFindAngle(A,O,B) 给出 $80$}
 &
 \tkzsubf{\begin{tikzpicture}
 \tkzDefPoint(0,0){O}    \tkzDefPoint(0:2){A}
 \tkzDefPointBy[rotation=center O angle -80](A)  \tkzGetPoint{B}
 \tkzDrawSegments[-{Stealth}](O,A O,B)
 \tkzMarkAngles[size=1,-Stealth,red](A,O,B)
 \tkzFindAngle(A,O,B)   \tkzGetAngle{an}
 \tkzLabelAngle[pos=.75,red](A,O,B){$\pgfmathprintnumber{\an}^\circ$}
\tkzLabelPoints[right](A)  \tkzLabelPoints[below](B)
 \end{tikzpicture}}
  {{\textbackslash}tkzFindAngle(A,O,B) 给出 $\pgfmathprintnumber{\an}^\circ$}
 \\\hline
 \end{tabular}
 \end{figure}

如我们所见，$-80^\circ$ 旋转定义了一个顺时针角，但宏命令
\tkzcname{tkzFindAngle} 恢复出的是一个逆时针角。

\subsection{获取角度 \tkzcname{tkzGetAngle}}
\begin{NewMacroBox}{tkzGetAngle}{\parg{name of macro}}%
将一个角的度数赋值给某个宏。该值为正数，介于 $0^\circ$ 与 $360^\circ$ 之间。本宏命令获取 \tkzcname{tkzAngleResult} 并存入一个新宏。

\medskip

\begin{tabular}{lll}%
\toprule
参数             & 示例 & 说明             \\
\midrule
\TAline{宏名称} {\tkzcname{tkzGetAngle}\{ang\}}{\tkzcname{ang} 中保存角度值。}
\end{tabular}
\end{NewMacroBox}

这个辅助宏获取下一个宏命令 \tkzcname{tkzFindAngle} 的结果。

\subsection{由三个点构成的角}

\begin{NewMacroBox}{tkzFindAngle}{\parg{pt1,pt2,pt3}}%
结果保存在宏 \tkzcname{tkzAngleResult} 中。

\medskip

\begin{tabular}{lll}%
\toprule
参数     & 示例 & 说明     \\
\midrule
\TAline{(pt1,pt2,pt3)} {\tkzcname{tkzFindAngle}(A,B,C)}{\tkzcname{tkzAngleResult} 给出角（$\overrightarrow{BA},\overrightarrow{BC}$）}
\bottomrule
\end{tabular}

\medskip
角度大小恒为正，介于 $0^\circ$ 与 $360^\circ$ 之间。按照通常约定，小于平角的逆时针角大小介于 $0^\circ$ 与 $180^\circ$ 之间，而小于平角的顺时针角大小大于 $180^\circ$。\tkzcname{tkzGetAngle} 用于获取该角。
\end{NewMacroBox}
 
\subsubsection{角度度量的验证}
    
\begin{tkzexample}[latex=7cm,small]
\begin{tikzpicture}[scale=.75]
  \tkzDefPoint(-1,1){A}
  \tkzDefPoint(5,2){B}
  \tkzDefEquilateral(A,B)
  \tkzGetPoint{C}
  \tkzDrawPolygon(A,B,C)
  \tkzFindAngle(B,A,C) \tkzGetAngle{angleBAC}
  \edef\angleBAC{\fpeval{round(\angleBAC)}}
  \tkzDrawPoints(A,B,C) 
  \tkzLabelPoints(A,B)
  \tkzLabelPoint[right](C){$C$}
  \tkzLabelAngle(B,A,C){\angleBAC$^\circ$}
  \tkzMarkAngle[size=1.5](B,A,C)
\end{tikzpicture}
\end{tkzexample}

\subsubsection{确定三角形的三个角}

\begin{tkzexample}[latex=6cm,small]
\begin{tikzpicture}
\tikzset{label angle style/.append style={pos=1.4}}
\tkzDefPoints{0/0/a,5/3/b,3/6/c}
\tkzDrawPolygon(a,b,c)
\tkzFindAngle(c,b,a)\tkzGetAngle{angleCBA}
\pgfmathparse{round(1+\angleCBA)}
\let\angleCBA\pgfmathresult
\tkzFindAngle(a,c,b)\tkzGetAngle{angleACB}
\pgfmathparse{round(\angleACB)}
\let\angleACB\pgfmathresult
\tkzFindAngle(b,a,c)\tkzGetAngle{angleBAC}
\pgfmathparse{round(\angleBAC)}
\let\angleBAC\pgfmathresult
\tkzMarkAngle(c,b,a)
\tkzLabelAngle(c,b,a){\tiny $\angleCBA^\circ$}
\tkzMarkAngle(a,c,b)
\tkzLabelAngle(a,c,b){\tiny $\angleACB^\circ$}
\tkzMarkAngle(b,a,c)
\tkzLabelAngle(b,a,c){\tiny $\angleBAC^\circ$}
\end{tikzpicture}
\end{tkzexample}

\subsubsection{两圆之间的夹角}
我们求两圆在交点处由切线所构成的角。

\begin{tkzexample}[latex=7cm,small]
\begin{tikzpicture}[scale=.4]
\pgfkeys{/pgf/number format/.cd,%
          fixed,precision=1}
\tkzDefPoints{0/0/A,6/0/B,4/2/C}
\tkzDrawCircles(A,C B,C)
\tkzDefLine[tangent at=C](A) \tkzGetPoint{a}
\tkzDefPointsBy[symmetry = center C](a){d}
\tkzDefLine[tangent at=C](B) \tkzGetPoint{b}
\tkzDrawLines[add=1 and 4](a,C  C,b)
\tkzFillAngle[fill=teal,opacity=.2%
                        ,size=2](b,C,d)
\tkzFindAngle(b,C,d)\tkzGetAngle{bcd}
\tkzLabelAngle[pos=1.25](b,C,d){%
  \tiny $\pgfmathprintnumber{\bcd}^\circ$}
\end{tikzpicture}
\end{tkzexample}

\subsection{直线与水平轴的夹角 \tkzcname{tkzFindSlopeAngle}}
比上一个更有趣：结果介于 $-180^\circ$ 与 $+180^\circ$ 之间。

\begin{NewMacroBox}{tkzFindSlopeAngle}{\parg{A,B}}%
确定直线 $(AB)$ 的斜率（以角度表示）。结果保存在宏 \tkzcname{tkzAngleResult} 中。

\medskip
\begin{tabular}{lll}%
\toprule
参数  & 示例 & 说明     \\
\midrule
\TAline{(pt1,pt2)} {\tkzcname{tkzFindSlopeAngle}(A,B)}{}
\bottomrule
\end{tabular}

\medskip
\tkzcname{tkzGetAngle} 用于获取结果；若不需要保留，可直接使用 \tkzcname{tkzAngleResult}。
\end{NewMacroBox}

\subsubsection{如何使用 \tkzcname{tkzFindSlopeAngle}}

 这里 $(AB)$ 是 $\widehat{CAD}$ 的角平分线，$(AD)$ 的斜率为零。我们先求 $(AB)$ 的斜率，然后旋转两次。

\begin{tkzexample}[latex=7cm,small]
\begin{tikzpicture}
 \tkzDefPoint(1,5){A} \tkzDefPoint(5,2){B}  
 \tkzFindSlopeAngle(A,B)\tkzGetAngle{tkzang}
 \tkzDefPointBy[rotation= center A angle \tkzang ](B)
 \tkzGetPoint{C}
 \tkzDefPointBy[rotation= center A angle -\tkzang ](B) 
 \tkzGetPoint{D}
 \tkzDrawSegment(A,B)
 \tkzDrawSegments[new](A,C A,D)
 \tkzDrawPoints(A,B,C,D)
 \tkzCompass[length=1](A,C)
 \tkzCompass[delta=10,brown](B,C)  
 \tkzLabelPoints(B,C,D)  
 \tkzLabelPoints[above left](A)
\end{tikzpicture}
\end{tkzexample}

\subsubsection{\tkzcname{tkzFindSlopeAngle} 与 \tkzcname{tkzGetAngle} 的使用}
这是另一种构造垂直平分线（中垂线）的方法。

\begin{tkzexample}[latex=6cm,small]
\begin{tikzpicture}
 \tkzInit
 \tkzDefPoint(0,0){A}        \tkzDefPoint(3,2){B}
 \tkzDefLine[mediator](A,B)  \tkzGetPoints{I}{J}
 \tkzCalcLength(A,B)         \tkzGetLength{dAB}
 \tkzFindSlopeAngle(A,B)     \tkzGetAngle{tkzangle}
 \begin{scope}[rotate=\tkzangle]
   \tkzSetUpArc[color=gray,line width=0.2pt,%
     /tkzcompass/delta=10]
   \tkzDrawArc[R,arc](B,3/4*\dAB)(120,240)
   \tkzDrawArc[R,arc](A,3/4*\dAB)(-45,60)
   \tkzDrawLine(I,J)         \tkzDrawSegment(A,B)
  \end{scope}
  \tkzDrawPoints(A,B,I,J)    \tkzLabelPoints(A,B)
   \tkzLabelPoints[right](I,J)
\end{tikzpicture}
\end{tkzexample}

\subsubsection{\tkzcname{tkzFindSlopeAngle} 的另一用法}

\begin{tkzexample}[latex=7cm,small]
\begin{tikzpicture}[scale=1.5]
  \tkzDefPoint(1,2){A}    \tkzDefPoint(3,4){B}
  \tkzDefPoint(3,2){C}    \tkzDefPoint(3,1){D}
  \tkzDrawSegments(A,B A,C A,D)
  \tkzDrawPoints[color=red](A,B,C,D)  
  \tkzLabelPoints(A,B,C,D)
  \tkzFindSlopeAngle(A,B)\tkzGetAngle{SAB} 
  \tkzFindSlopeAngle(A,C)\tkzGetAngle{SAC}
  \tkzFindSlopeAngle(A,D)\tkzGetAngle{SAD}
  \pgfkeys{/pgf/number format/.cd,fixed,precision=2}
  \tkzText(1,5){The slope of (AB) is : 
     $\pgfmathprintnumber{\SAB}^\circ$}     
  \tkzText(1,4.5){The slope of (AC) is : 
     $\pgfmathprintnumber{\SAC}^\circ$}    
  \tkzText(1,4){The slope of (AD) is : 
     $\pgfmathprintnumber{\SAD}^\circ$}
\end{tikzpicture}
\end{tkzexample}

\endinput
